% answers_short.tex —— 2009 年·普通高等学校招生考试·广东卷·理科 · 简短答案速查
% 由 tools/make_short_answers.py 自 answers.tex 的 \gksolhead/\gkanswer 提取，
% 并与 2009/2009_数学真题_广东理_答案要点.md 逐题比对后生成，勿手改（改源 MD 后重跑）。

\begin{gknotice}
\textbf{使用说明}：本节只给最终答案便于快速核对，完整推导见解析版。答案为编写者独立推导，\textbf{非官方标准答案}。
\end{gknotice}

\gkpart{选择题}
\begin{center}\begin{tabular}{ccccccccc}
\toprule
题号 & 1 & 2 & 3 & 4 & 5 & 6 & 7 & 8 \\
\midrule
答案 & B & C & B & C & D & D & A & A \\
\bottomrule
\end{tabular}\end{center}

\gkpart{填空题}
\begin{itemize}[leftmargin=2.6em]
  \item[\textbf{9.}] $s=\dfrac{a_1+a_2+\cdots+a_n}{n}$，样本的平均数。
  \item[\textbf{10.}] $\boldsymbol a=(-1,1)$ 或 $\boldsymbol a=(-3,1)$。
  \item[\textbf{11.}] $\dfrac{x^2}{36}+\dfrac{y^2}{9}=1$。
  \item[\textbf{12.}] $a=\dfrac{5}{12}$，$b=\dfrac14$。
  \item[\textbf{13.}] $-1$。
  \item[\textbf{14.}] $(-\infty,-2)\cup\left(-2,-\dfrac32\right]$。
  \item[\textbf{15.}] $8\pi$。
\end{itemize}

\gkpart{解答题}
\begin{itemize}[leftmargin=2.6em]
  \item[\textbf{16.}] （1）$\dfrac25$；（2）$\dfrac{\sqrt2}{2}$。
  \item[\textbf{17.}] （1）$x=\dfrac{119}{18250}$；（2）良 $119$ 天，轻微污染 $100$ 天；（3）$\dfrac{76653}{78125}$。
  \item[\textbf{18.}] （1）$\dfrac23$；（2）证明见解析；（3）$\dfrac{\sqrt3}{3}$。
  \item[\textbf{19.}] （1）$y=2x^2-x+\dfrac{11}{8}$（$-\dfrac14<x<\dfrac54$）；（2）$-\dfrac{7\sqrt2}{5}$。
  \item[\textbf{20.}] （1）$m=\sqrt2-1$ 或 $m=-1-\sqrt2$；（2）分类见解析。
  \item[\textbf{21.}] （1）$x_n=\dfrac{n}{n+1}$，$y_n=\dfrac{n\sqrt{2n+1}}{n+1}$；（2）证明见解析。
\end{itemize}
